\subsection{Nullstellen und Definitionslücken}
\label{subsec:real_functions_poles}
\begin{DefinitionBox}{Nullstellen und Definitionslücken}{real_functions_rational_zeros_and_gaps}
  \label{def:real_functions_removable_gap}
  \label{def:real_functions_pole}
  Die Nullstellen von Zähler $Z$ und Nenner $N$ erlauben folgende Einteilung:
  \begin{center}
    \CachedTikzFigure{content/resources/chapters/real_functions/tikz/functions_rational_case_hierarchy.tex}
  \end{center}
  Das folgende Ablaufdiagramm zeigt die Unterscheidung im Einzelnen.
  \begin{center}
    \CachedTikzFigure{content/resources/chapters/real_functions/tikz/functions_rational_decision.tex}
  \end{center}
  Polstellen und hebbare Definitionslücken sind beides Definitionslücken.
\end{DefinitionBox}

\begin{ExampleBox}{Gebrochen-rationale Funktion mit Nullstelle}{real_functions_rational_zero_example}
  Wir untersuchen die Funktion $f(x)$ auf Nullstellen und Definitionslücken.
  \[
    f(x)=\frac{x}{x^2+1}.
  \]
  Aus der Linearfaktorzerlegung von Zähler und Nenner lesen wir deren Nullstellen ab:
  \begin{ClozeArea}[32]{6}
    \[
      \begin{alignedat}{3}
        Z(x)&=x
          &\quad\Longleftrightarrow\quad& Z(x)=(x-0)
          &\quad\Longrightarrow\quad& \text{Nullstelle bei }x_1=0,\\[0.8em]
        N(x)&=x^2+1
          &\quad\Longleftrightarrow\quad& N(x)=x^2+1\ge1
          &\quad\Longrightarrow\quad& \textcolor{red}{\text{keine Nullstelle}}.
      \end{alignedat}
    \]
  \end{ClozeArea}
  \noindent
  \begin{minipage}[c]{0.45\linewidth}
    Da der Nenner keine Nullstellen hat, liegt keine Definitionslücke vor:
    \[
      D_f=\Cloze{2cm}{\R}.
    \]
    Die Zählernullstelle $x_1=0$ liegt im Definitionsbereich $D_f$ und ist damit
    eine Nullstelle von $f$:
    \[
      f(0)=\Cloze{0.8cm}{0}.
    \]
  \end{minipage}\hfill
  \begin{minipage}[c]{0.52\linewidth}
    \centering
    \CachedTikzFigure{content/resources/chapters/real_functions/tikz/functions_rational_zero.tex}
  \end{minipage}
\end{ExampleBox}

\begin{ExampleBox}{Gebrochen-rationale Funktion mit Polstelle}{real_functions_rational_pole_example}
  Wir untersuchen die Funktion $g(x)$ auf Nullstellen und Definitionslücken.
  \[
    g(x)=\frac{1}{x-1}.
  \]
  Aus der Linearfaktorzerlegung von Zähler und Nenner lesen wir deren Nullstellen ab:
  \begin{ClozeArea}[32]{6}
    \[
      \begin{alignedat}{3}
        Z(x)&=1
          &\quad\Longleftrightarrow\quad& Z(x)=1
          &\quad\Longrightarrow\quad& \textcolor{red}{\text{keine Nullstelle}},\\[0.8em]
        N(x)&=x-1
          &\quad\Longleftrightarrow\quad& N(x)=(x-1)
          &\quad\Longrightarrow\quad& \text{Nullstelle bei }x_1=1.
      \end{alignedat}
    \]
  \end{ClozeArea}
  \noindent
  \begin{minipage}[c]{0.45\linewidth}
    Die Nennernullstelle $x_1=1$ ist eine Definitionslücke:
    \[
      D_g=\Cloze{3cm}{\R\setminus\{1\}}.
    \]
    Da aber $Z(1)=1 \neq 0$, liegt bei $x=1$ eine Polstelle vor.
    Zu jeder Polstelle gehört eine \textcolor{green!55!black}{\emph{vertikale Asymptote}}
    \EnglishTerm{vertical asymptote} an dieser Stelle.
  \end{minipage}\hfill
  \begin{minipage}[c]{0.52\linewidth}
    \centering
    \CachedTikzFigure{content/resources/chapters/real_functions/tikz/functions_rational_pole.tex}
  \end{minipage}
\end{ExampleBox}

\begin{ExampleBox}{Gebrochen-rationale Funktion mit hebbarer Definitionslücke}{real_functions_cancelled_hole}
  Wir untersuchen die Funktion $f(x)$ auf Nullstellen und Definitionslücken.
  \[
    f(x)=\frac{x^2-1}{x+1}.
  \]
  Aus der Linearfaktorzerlegung von Zähler und Nenner lesen wir deren Nullstellen ab:
  \begin{ClozeArea}[34]{5}
    \[
      \begin{alignedat}{3}
        Z(x)&=x^2-1
          &\quad\Longleftrightarrow\quad& Z(x)=(x-1)(x+1)
          &\quad\Longrightarrow\quad& \text{Nullstellen bei }x_1=-1,\;x_2=1,\\[0.8em]
        N(x)&=x+1
          &\quad\Longleftrightarrow\quad& N(x)=(x+1)
          &\quad\Longrightarrow\quad& \text{Nullstelle bei }x_1=-1.
      \end{alignedat}
    \]
  \end{ClozeArea}
  \noindent
  \begin{minipage}[c]{0.45\linewidth}
    Die Nennernullstelle $x_1=-1$ ist eine Definitionslücke:
    \[
      D_f=\Cloze{3cm}{\R\setminus\{-1\}}.
    \]
    Die Zählernullstelle $x_2=1$ liegt im Definitionsbereich $D_f$ und ist damit
    eine Nullstelle von $f$:
    \[
      f(x_2)=f(1)=\Cloze{0.8cm}{0}.
    \]
  \end{minipage}\hfill
  \begin{minipage}[c]{0.52\linewidth}
    \centering
    \CachedTikzFigure{content/resources/chapters/real_functions/tikz/functions_removable_gap.tex}
  \end{minipage}
  \par\medskip
  Bei $x_1=-1$ können wir $f$ nicht auswerten.
  Der Graph zeigt jedoch, welchem Wert sich die Funktion dort nähert.
  Diesen Grenzwert bestimmen wir, indem wir in der Linearfaktorzerlegung von
  Zähler und Nenner den gemeinsamen Faktor $(x+1)$ kürzen:
  \begin{ClozeArea}[34]{6}
    \[
      f(x)=\frac{(x-1)\cancel{(x+1)}}{\cancel{(x+1)}}=x-1\qquad(x\ne-1).
    \]
  \end{ClozeArea}
  In den Term $x-1$ lässt sich auch $x=-1$ einsetzen.
  Deshalb definieren wir mit demselben Term die Funktion $\widetilde f$ auf ganz $\R$:
  \[
    \widetilde f(x)=x-1 \qquad(x\in\R).
  \]
  Die Funktion $\widetilde f$ ist eine \emph{stetige Fortsetzung}
  \EnglishTerm{continuous extension} von $f$.
  Sie unterscheidet sich von $f$ nur dadurch, dass sie auch bei $x=-1$ definiert ist.

  Den gesuchten Grenzwert können wir daher direkt mit der stetigen Fortsetzung bestimmen:
  \[
    \lim_{x\to-1}f(x)=\widetilde f(-1)=-2.
  \]
\end{ExampleBox}
